\section{Conclusion}
\label{Section:Conclusion}

This work set out to determine whether a single deep reinforcement learning design, trained separately for each instrument and otherwise unchanged, produces a profitable and risk-controlled policy on all seven major currency pairs when the evaluation charges what a real account charges.

The answer is a qualified yes, and the qualification is as much a part of the result as the answer.

\subsection{Summary}
\label{Subsection:Summary}

Three things were built and one question was answered.

The first is a trading framework in which one strategy implementation runs in three places: an offline backtest over stored quotes, the cTrader backtesting tab, and a live account. Its execution engine reproduces the platform's own results byte for byte across six reference configurations and extends them, stepping on ticks rather than bar closes and charging the spread recorded at the moment of each trade. Behind it sits a market database of $1.66$ billion quotes across the seven pairs, occupying $295$ GB, from which every bar series used in this work is derived.

The second is an agent. It observes thirty-eight dimensionless values built from sixteen technical indicators and emits one continuous action that sets the direction and the size of its position together. Six deviations from the published algorithm were needed to make it train on this problem, of which two matter: a scale-sensitive reward, because a scale-invariant one gives no gradient towards holding a meaningful position, and a penalty on the actor's pre-activation, without which the policy saturates and collapses to a constant action.

The third is a selection procedure. Candidates are rejected by gates that remove degenerate policies, tested against a null built by rotating each model's own position sequence in time, and ranked by a composite in which return is one component of nine. The procedure was not built to be elegant; it was built because ranking on return does not work, which Section~\ref{Subsection:StatisticalScreening} demonstrates rather than asserts.

The question was answered as follows. All seven agents are profitable over the eleven-year window and from all five starting balances, with total returns from $+15.7\%$ to $+226.6\%$ net of the spread and commission of a swap-free raw-spread account. Six of the seven produce positive alpha once their return is decomposed against their own market, and three of those carry almost no market exposure while doing it. The seventh, USD/JPY, reproduces its market and is reported as a failure with a diagnosed cause. Combining the seven into an equal-weight portfolio reduces the maximum drawdown from an average of $34.5\%$ to $15.9\%$ and raises the Sharpe ratio to $0.632$, because the agents are close to uncorrelated with one another.

Two findings sit underneath those numbers and are more transferable than they are. The continuous action is genuinely used: the mean absolute action ranges from $0.09$ to $0.98$ across the pairs and saturation is below one percent everywhere, so the agents learned to size positions rather than to switch between extremes. And win rates lie between $49.8\%$ and $55.0\%$, with the strongest model in the set having the lowest of them, so the edge does not come from being right more often. It comes from what the agents do with a position once it is open, which is the capability a continuous action provides and a discrete one does not.

\subsection{Limitations}
\label{Subsection:Limitations}

The most important limitation is stated first because it bounds everything else. The performance figures reported in Section~\ref{Section:ExperimentalResults} are measured over a window that includes the years the agents trained on. The procedure does hold a year back from training and from the choice among an agent's own episodes, but the choice of which trained agent to publish was made on the full window, so the headline returns are an in-sample result. On the one year that no part of the procedure saw, the models lost money: the portfolio returned $-9.8\%$ in 2025 and five of the seven pairs were negative. That year contained a broad dollar reversal the training decade did not, and the one model that held up is the least directional of the set, but a single year is a small sample and nothing stronger should be read into it in either direction.

Selection is uncorrected in a second sense. Each published model was chosen from between $128$ and $256$ candidates, so a permutation $p$-value near $0.05$ is close to what that many draws produce by chance. The screening is therefore reported as a filter that demonstrably rejected bad models, not as a significance claim for the ones that survived. Under a strict correction only USD/CAD and GBP/USD would clearly stand.

Three narrower limitations belong on the record. Position size was chosen per pair after results were seen; it scales every trade without changing the policy that generates them, but it is a parameter fitted on the same data as everything else. Training is not reproducible across environments: it is bit-exact within one installation at a single thread, and a change of numerical library version alters the outcome, so the delivered models are archived artifacts with a verified replay rather than a recipe that regenerates them. And two weeks of January 2016 are missing from all seven pairs for reasons on the broker's side; the gap is uniform across pairs and falls outside the evaluation window, so it does not bias any comparison made here.

Finally, the scope is narrower than the method. Seven instruments in one asset class over one eleven-year period, with one decision cadence and one cost model, is a single point in a large space. Nothing in this work establishes that the design transfers to other markets, other periods or other holding horizons, and the USD/JPY result shows that it does not transfer to every instrument even within the seven.

\subsection{Future Work}
\label{Subsection:FutureWork}

The limitations above order the work that follows, because the most useful next step is the one that would most change what the results are worth.

The first is a genuinely held-out evaluation. Training would end before the reporting period begins, and the choice among candidates would be made using only data from before that boundary, so that the reported figures are an out-of-sample estimate rather than a description of the window the models were selected on. This is the single change that would most strengthen the claims, and the framework already supports it; what it needs is a campaign run under the stricter protocol rather than new machinery.

The second follows from it. The procedure used here divides the window once. A walk-forward scheme with several folds, in which each fold trains on the data before it and is evaluated on the data after, would produce an out-of-sample record for the whole period rather than for a single year. Carrying the learned weights from one fold into the next, rather than starting each fold from a fresh initialisation, would test whether an agent adapts as the market changes; this is supported by the framework and was not exercised in this campaign.

The third is to report the variation between training runs. The path-robustness protocol used here varies the starting balance, which measures sensitivity to the entry path but not to the training itself. Reporting the distribution of outcomes across the candidate pool, and where the published model sits within it, would let a reader see how much of a result is method and how much is draw.

Beyond the evaluation, three directions follow from what the agents did. The portfolio of Section~\ref{Subsection:Portfolio} weights the seven equally; weighting by risk, or by the correlation structure, should improve on that, and the near-independence of the components suggests the gain would be real. The observation is a single frame of thirty-eight values, so a recurrent or attention-based encoder that reads a window of history directly is the natural extension, and it was in the original plan for this work. And the framework already runs the same strategy on a live account, so a sustained live deployment would test the one thing no backtest can, which is whether the execution assumptions hold when the orders are real.

The USD/JPY result deserves its own line. The failure there is understood as a consequence of a monotonic price window offering no gradient towards modulating exposure. Whether that is the whole explanation is not established, and the pair should be revisited once the evaluation protocol above is in place.

Deep reinforcement learning applied to currency trading is neither the solved problem that some reported returns suggest nor the dead end that its out-of-sample record would imply. What this work shows is that the difficulty is not mainly in the learning algorithm. It is in the environment the agent is trained against, the costs it is charged, and above all in the discipline of deciding which of many trained models deserves to be believed. Six of seven agents produced returns that survive a decomposition into skill and drift; in every one of the seven the candidate that looked best on return did not.